\hypertarget{appendix-41e8ef22a21e}{}
\section*{Fused multiply-add}
\section{Fused multiply-add}\label{fused-multiply-add}

\textbf{Curriculum depth:} Introduced · \textbf{First seen:} C005 · \textbf{Latest development:} C005

\subsection{Definition and mechanism}\label{definition-and-mechanism}

FMA rounds ab+c once rather than rounding the product first. For p=2\^{}-23,a=b=1+p,c=−(1+2p), separate binary32 gives zero, FMA gives exact 2\^{}-46. Compiler contraction can change an unfused-looking expression.

\subsection{Prerequisites}\label{prerequisites}

\begin{itemize}
\tightlist
\item
  \hyperref[appendix-343615f49ddf]{IEEE-754 floating point} --- Fused operation semantics and destination rounding are part of the arithmetic contract.
\item
  \hyperref[appendix-614afe1de592]{Rounding to nearest ties to even} --- Midpoint decisions determine separate/fused results and narrowed persistent state.
\end{itemize}

\subsection{Used by}\label{used-by}

\begin{itemize}
\tightlist
\item
  \hyperref[appendix-c2ac4a84c283]{Floating-point reassociation} --- uses this mechanism to establish its numerical contract.
\item
  \hyperref[appendix-5aef4f13483a]{Dot and cross-product rounding} --- uses this mechanism to establish its numerical contract.
\end{itemize}

\subsection{Recall and next depth}\label{recall-and-next-depth}

Reconstruct the definition's example and name the assumptions that make it valid.

\textbf{Pending:} Specify FMA policy for reference/reproducible builds; a fused chain is not a fully exact dot product.

\subsection{Source and navigation}\label{source-and-navigation}

\href{Continuous\%20Notes/Computing\%20Foundations\%20-\%20From\%20Bits\%20to\%20C++\%20and\%20Systems\%20APIs\#When\%20arithmetic\%20meets\%20uncertainty}{M005 teaching chapter} · \hyperref[appendix-fd090d982a15]{Part definitions} · \href{}{Dependency Map (02-Foundations-II, appendix-4ad47bdd269a)}
